\subsection*{Appendix A.5: Heuristic Strategy for Column Pair Selection}
\label{app:column_selection}

In the R-step, selecting the most effective column pairs for rotation is critical for efficiency. We propose a greedy strategy based on \textbf{Imbalance} and \textbf{Complementarity} to prioritize pairs that yield the maximum potential for distribution reshaping.

\textbf{1. Column Imbalance Score}

Given the normalized activation matrix $\mathbf{Z} = \mathbf{S}^{-1}\mathbf{X}\mathbf{R}$, let $\mathbf{z}^{(k)}$ denote the $k$-th column vector. We first construct the codebook histogram for each column:
\begin{equation}
N_j^{(k)} = \sum_{b=1}^{B} \mathbb{1}\left( |z_{b,k}| \in \mathcal{I}_j \right), \quad \hat{p}_j^{(k)} = \frac{N_j^{(k)}}{\sum_{l=1}^J N_l^{(k)}}
\end{equation}
We define the \textbf{Imbalance Score} $h_k$ as the mean squared distance between the empirical distribution and the uniform distribution:
\begin{equation}
h_k = \sum_{j=1}^{J} \left( \hat{p}_j^{(k)} - \frac{1}{J} \right)^2
\end{equation}
A larger $h_k$ indicates that the column's activations are highly concentrated in a few codewords (e.g., collapsed to 0 or max), making it a priority target for optimization. We sort all columns by $h_k$ in descending order to form a priority queue.

\textbf{2. Pair Complementarity}

Simply selecting the two columns with the highest $h_k$ is suboptimal if they share similar skewness (e.g., both biased towards large values). We introduce \textbf{Complementarity} to measure the statistical difference between two columns.
We define the "anti-correlation" $c_{k,l}$ between column $k$ and column $l$ as:
\begin{equation}
c_{k,l} = - \sum_{j=1}^{J} \left( \hat{p}_j^{(k)} - \frac{1}{J} \right) \left( \hat{p}_j^{(l)} - \frac{1}{J} \right)
\end{equation}
A positive $c_{k,l}$ implies that when one column over-occupies a specific codeword, the other tends to under-occupy it. This "complementary" nature facilitates the efficient transfer of probability mass between the two columns via 2D rotation.

Combining both metrics, we define the \textbf{Pair Selection Score} $H_{k,l}$:
\begin{equation}
H_{k,l} = h_k + h_l + \lambda c_{k,l}
\end{equation}
where $\lambda > 0$ balances individual urgency and mutual compatibility.

\textbf{3. Selection Workflow}

In each R-step iteration, the selection proceeds as follows:
\begin{enumerate}
    \item \textbf{Filtering:} Calculate $h_k$ for all columns and select the top $K_{\text{top}}$ (e.g., $K/2$) most imbalanced columns as the candidate pool.
    \item \textbf{Pairing:} Compute $H_{k,l}$ for all pairs within the candidate pool.
    \item \textbf{Execution:} Select the top $P$ non-overlapping pairs with the highest $H_{k,l}$ scores.
    \item \textbf{Optimization:} Apply the exact angle solver (Appendix A.4) to these $P$ pairs to update $\mathbf{R}_{\text{intra}}$.
\end{enumerate}
This strategy ensures that limited computational resources are focused on the "most problematic" columns and the "most fixable" pairs.